% MODELO COBEF 2021 adaptado AC Araujo
\documentclass[10pt,fleqn,a4paper,twoside]{article}
\usepackage{abcm}
\usepackage[export]{adjustbox}

\usepackage[figuresright]{rotating}
\usepackage{amsmath,amssymb}
\usepackage{hyperref}
\usepackage{subfigure}
\usepackage{notoccite}
\usepackage{amsmath,amssymb}
\usepackage{hyperref}
\usepackage{subfigure}
\usepackage[brazil]{babel}
\usepackage[utf8]{inputenc}
\usepackage[T1]{fontenc}
\usepackage{graphicx}
\usepackage{amsmath,amssymb}
\usepackage{algorithmic}
\usepackage{algorithm}
\usepackage{float}
\usepackage{color}
\usepackage{epstopdf}
\usepackage{epsfig}
\usepackage{pdfpages}
\usepackage[export]{adjustbox}
\usepackage[section]{placeins}


\usepackage{bigstrut}
\usepackage{rotating}
\usepackage{multirow}
\usepackage{array}
\usepackage{dcolumn}
\usepackage{booktabs}

% Pacote para melhorar opções de tabelas
\usepackage[x11names,table]{xcolor}
\usepackage{float}
\usepackage{threeparttable}
\usepackage{scalefnt}

\makeatletter
\setlength{\@fptop}{0pt}
\makeatother	

%COBEF
\pagestyle{fancy}
\fancyfoot{}
\setlength{\headheight}{30pt}
\fancyhead[RO]{11º Congresso Brasileiro de Engenharia de Fabricação \\ 24 a 26 de maio de 2021, Curitiba, PR, Brasil}
\fancyhead[LE]{Autor, S. Autor e T. Autor \\ Titulo do trabalho em uma linha}
\fancyhead[Lo,RE]{}


\begin{document}
    
% CABEÇALHO
\thispagestyle{empty}
\begin{figure}[ht]
\begin{center}
 \includegraphics[scale=0.5]{cabecalho2021.jpg}
\end{center}
\end{figure}
\vspace{0.5cm}
 \hspace{-.8cm}

\begin{tabular}{||p{\textwidth}}
\begin{center}
\vspace{-.6cm}
\title{INSTRUÇÕES PARA FORMATAÇÃO DOS ANAIS DO COBEF 2021}
\end{center}
\authors{Primeiro Autor $^1$} \\
\authors{Segundo Autor $^2$} \\
\institution{$^{1,2}$ Instituicao 1} \\
\email{$^1$ email1@mecanica.uf.br, $^2$ email2@materiais.ue.br}
	\\
	\\
\authors{Autor3$^3$} \\
\institution{$^3$ Instituicao 2} \\
\email{email3@producao.cefet.br}
    \\
	\\
    \abstract{\textbf{Resumo.} O objetivo dessas instruções é servir como um guia para a formatação dos artigos a serem publicados nos anais do 11$^o$ Congresso Brasileiro de Engenharia de Fabricação, COBEF 2021. O resumo deve descrever os objetivos, metodologia e conclusões principais do trabalho em menos de 200 palavras. Ele não deve conter equações ou referências bibliográficas.  }\\
    \keywords{\textbf{Palavras-chave:} Palavra 1. Palavra 2. Palavra 3 … (até 5 palavras) }\\
    \end{tabular}
    
 \section{INTRODUÇÃO}

Os artigos devem ser formatados estritamente de acordo com as instruções formatadas no template Word do COBEF, que inspirou este template. 
Os artigos devem ter no máximo 8 páginas, incluindo as seguintes seções: título, afiliação, resumo, palavras-chave, introdução, metodologia, resultados e conclusões. Sua versão final deverá ser submetida em formato PDF e não deverá exceder 2.0 Mb.



\section{FORMATO DO TEXTO}

O texto deverá ser redigido em português ou inglês em páginas de tamanho A4, usando fonte Times New Roman, tamanho 10, exceto para o título, identificação de autores, afiliação, resumo e palavras chave, que devem seguir as formatações indicadas acima. A primeira página deve ter margem superior de 3 cm e as demais margens de 2 cm. Todas as margens das demais páginas devem ter 2 cm.

AS PÁGINAS NÃO DEVEM SER NUMERADAS.

O bloco de texto contendo título, identificação de autores, afiliação, resumo e palavras chave deve ser recuado 0,1 cm da margem esquerda e destacado por uma barra vertical de espessura 2 ¼ pt na borda esquerda.

O texto deve ter alinhamento justificado, tendo a primeira linha de cada parágrafo um recuo de 0,5 cm. Notas de rodapé devem ser evitadas.

Símbolos e notações devem ser descritos no texto e as grandezas físicas expressas no sistema internacional. Símbolos matemáticos devem ser digitados em itálico.

No word: As referências bibliográficas devem ser citadas no texto pelo último nome dos autores e o ano de publicação, de acordo com os seguintes exemplos: “Trabalhos recentes (de Oliveira e Melo, 2013)...” ou “Recentemente, McConnell e Varoto (2008)...”. .... No caso de haver três ou mais autores, a forma “Rodriguez, et al. (2006)” deve ser utilizada. Duas ou mais referências com os mesmos autores e anos de publicação devem ser diferenciadas pelos índices “a”, “b”, etc. após o ano de publicação. Por exemplo: “Trabalho recente (Trindade e Benjeddou, 2011a)...”.

No Latex: As referências bibliográficas devem ser citadas no texto pelo último nome dos autores e o ano de publicação, de acordo com os seguintes exemplos: “Trabalhos recentes \citep{Ref1}...” ou “Recentemente, \cite{Ref1}....”. No caso de haver três ou mais autores, a forma \cite{livro1} deve ser utilizada. 
Duas ou mais referências com os mesmos autores e anos de publicação devem ser diferenciadas pelos índices “a”, “b”, etc. após o ano de publicação. Por exemplo: “Trabalho recente \citep{Ref2}...”.
 \nocite{Ref3}

Referências aceitáveis incluem artigos de revistas técnicas, dissertações, teses, anais de conferências, livros e comunicações pessoais. Páginas de internet também podem ser utilizadas.
As referências devem ser listadas ao final do trabalho, conforme instruções indicadas na Seção 4.

\subsection{Títulos e subtítulos de seções}

Os títulos e subtítulos de seções devem ser alinhados à esquerda e digitados com fonte Times New Roman, tamanho 10, em negrito. Eles devem ser numerados por meio de algarismos arábicos separados por pontos. Não mais do que 3 sub níveis devem ser utilizados e uma linha em branco deve ser inserida acima e abaixo de cada título/subtítulo. 

\subsection{Equações}

As equações devem ter recuo de 0,5 cm a partir da margem esquerda. Elas devem ser escritas com fonte Times New Roman, em itálico, com tamanho 10. Algarismos arábicos entre parênteses e alinhados à direita devem ser usados para a identificação das equações. No texto, as equações devem ser referenciadas como “Eq. (\ref{Eq1} )” no meio da frase e como “Equação (\ref{Eq1})” no início da frase.  Símbolos usados nas equações devem ser definidos imediatamente antes ou depois de sua primeira aparição.

Uma linha em branco deve ser inserida acima e abaixo de cada equação.

\begin{equation}
d\vec{F}=\vec{k} \: h \: db
\label{Eq1}
\end{equation}

\subsection{Figuras e tabelas}


Figuras e tabelas devem ser posicionadas o mais próximo possível de sua primeira citação e devem ser identificadas sequencialmente em numerais arábicos. Figuras e Tabelas devem ser referenciadas como “Fig. \ref{Fig1}” e “Tab. \ref{Tab1}” no meio da frase e como “Figura \ref{Fig1}” e “Tabela \ref{Tab1}” no início da frase. As figuras, tabelas e suas legendas devem ser centralizadas na página. As legendas, digitadas com Times New Roman tamanho 10, não devem ter mais do que 3 linhas.

Uma linha em branco deve ser inserida acima e abaixo de cada figura ou tabela. 

 \begin{figure}[ht]
\begin{center}
\includegraphics[scale=0.4]{figures/boxplotlongy.jpg}
\end{center}
\caption{Exemplo de figura}\label{Fig1}
\end{figure}

\begin{table}[htbp]
  \centering
  \caption{O estilo das bordas das tabelas é livre}
    \begin{tabular}{llllcc}
	\\
 \toprule
    \textbf{Fatores} & \multicolumn{3}{c}{\textbf{Níveis}} & \multicolumn{2}{c}{\textbf{Saídas}} \\
\midrule
    Longitude & Norte(+1) & Centro(0) & Sul(-1) & $S_x$(\%) & $S_y$(\%) \\
    Latitude & Oeste(+1) & Centro(0) & Leste(-1) & $S_x$(\%) & $S_y$(\%) \\
  \bottomrule
    \end{tabular}%
  \label{Tab1}%
\end{table}%





\section{CONCLUSÕES}

Exemplo, ..., Esse artigo abordou a influência que diferentes localizações na plataforma poderiam afetar as dimensões $x$ e $y$ de uma peça fabricada por SLA. 

Exemplo, ..., A partir da análise realizada, pode-se concluir que: 
\begin{itemize}
\item o fator latitude foi o único que apresentou influência sobre a contração na dimensão $x$;
\item o fator longitude foi o único que apresentou influência sobre a contração na dimensão $y$;

\end{itemize}

Exemplo, ..., Contudo, deve-se considerar os resultados experimentais ao projetar a fabricação por SLA. Todos os valores de contração encontrados foram menores que 5\% para um elemento de 3 mm de comprimento. 

\section{AGRADECIMENTOS}

Esta seção, opcional, deve ser posicionada antes da lista de referências.

\nocite{ABCM}
\nocite{phdthesis}
\section{REFERÊNCIAS}
\bibliographystyle{abcm}
\bibliography{bibliografia}

\section{RESPONSABILIDADE PELAS INFORMAÇÕES}

Os autores são os únicos responsáveis pelas informações incluídas neste trabalho.

      
\newpage

\noindent{
   \\ 
    \begin{tabular}{||p{\textwidth}}
    \begin{center}
    \vspace{-.6cm}
    \title{AFTER FULL PAPER IN PORTUGUESE OR SPANISH, IT’S NECESSARY THE ABSTRACT IN ENGLISH}
    \end{center}
    \authors{First Author’s Name $^1$} \\
    \authors{Second Author’s Name $^1$} \\
	 \institution{$^{1,2}$ Institution and address for first author} \\
\email{$^1$ email1@mecanica.uf.br, $^2$ email2@materiais.ue.br}
	\\
	\\
    \authors{Third Author’s Name $^3$} \\\\
    \institution{$^3$ Institution and address for third author} \\
\email{$^3$ email3@mecanica.uf.br}
	\\
    \\
    \abstract{\textbf{Resumo:} The purpose of these instructions is to serve as a guide for formatting papers to be published in the Proceedings}\\
    \keywords{\textbf{Palavras-chave:} keyword 1, keyword 2, keyword 3 (up to 5 keywords) }\\
    \end{tabular}
}

\end{document}